
\begin{frame}
  \titlepage
\end{frame}

\begin{frame}{真实长文档任务: 从冗长文本中找回关键事实}
  \centering
  \begin{minipage}{0.92\textwidth}
  \small
  \begin{columns}[T,onlytextwidth]
    \begin{column}{0.61\linewidth}
      \begin{block}{Context}
        {\fontsize{5.4pt}{6.5pt}\selectfont\ttfamily
        SUBSTANTIAL EQUIVALENCE DETERMINATION DECISION SUMMARY\\
        A. 510(k) Number: K143329\\
        B. Purpose for Submission: To obtain clearance for a new device\\
        C. Measurand: A conserved multi-copy sequence of Trichomonas vaginalis genomic DNA\\
        \textcolor{scured}{\textbf{D. Type of Test: Nucleic acid amplification assay}}\\
        \textcolor{scured}{\textbf{\quad (Helicase-dependent Amplification, HDA)}}\\
        E. Applicant: Quidel Corporation\\
        F. Proprietary and Established Names: Amplivue Trichomonas Assay\\
        G. Regulatory Information: Classification: Class II; Panel: Microbiology\\
        H. Intended Use: qualitative detection of Trichomonas vaginalis nucleic acids\\
        \hfill\textcolor{gray}{\ldots}
        }
      \end{block}
      \vskip0.18cm
      {\footnotesize\color{gray}真实文档中只有少数关键信息与问题相关, 其余内容都构成检索干扰.}
    \end{column}

    \begin{column}{0.35\linewidth}
      \vskip0.08cm
      长上下文场景中, 关键信息往往埋在很长、很密的文本中。

      \vskip0.35cm
      \begin{block}{Query}
        \centering
        \large What is the Type of Test?
      \end{block}

      \vskip0.20cm
      \begin{block}{Answer}
        \centering
        {\bfseries\color{scured}Nucleic acid amplification assay}\\
        {\scriptsize\color{scured}(Helicase-dependent Amplification, HDA)}
      \end{block}
    \end{column}
  \end{columns}

  \vskip0.34cm
  \begin{beamercolorbox}[sep=5pt,rounded=true,wd=\linewidth]{block body}
    \centering\small
    模型不仅要读完整段上下文, 还要在后续查询时准确定位并读出对应事实.
  \end{beamercolorbox}
  \end{minipage}
\end{frame}

\begin{frame}{长上下文模型真的“记住”了吗?}
  \centering
  \begin{minipage}{0.92\textwidth}
  \small
  长文档问答、代码理解、长对话都依赖前文信息。长上下文能力的关键不只是
  \textcolor{scured}{“能输入多长”}, 更重要的是: 后面能不能准确找回前面出现过的信息。

  \vskip0.32cm
  \begin{beamercolorbox}[sep=7pt,rounded=true,shadow=true,wd=\linewidth]{block body}
    \centering
    {\footnotesize\color{gray}Input: 前文写入若干键值绑定, 后文只给出查询键}\\[3pt]
    {\resizebox{0.94\linewidth}{!}{\ttfamily
      A\,4 \quad B\,3 \quad C\,6 \quad F\,1 \quad E\,2
      \quad \textcolor{scured}{\(\longrightarrow\)}
      \quad A\,? \quad C\,? \quad F\,? \quad E\,? \quad B\,?
    }}\\[5pt]
    {\footnotesize\color{gray}Output: 模型必须按查询顺序读出对应值}\\[3pt]
    {\Large\bfseries\color{scured}4,\;6,\;1,\;2,\;3}
  \end{beamercolorbox}

  \vskip0.30cm
  \begin{columns}[T,onlytextwidth]
    \begin{column}{0.48\linewidth}
      \begin{block}{形式上}
        \centering
        \vspace{0.12cm}
        记住前文的键值绑定
        \vspace{0.12cm}
      \end{block}
    \end{column}
    \begin{column}{0.48\linewidth}
      \begin{block}{本质上}
        \centering
        \vspace{0.12cm}
        在大量干扰信息中准确读出
        \vspace{0.12cm}
      \end{block}
    \end{column}
  \end{columns}

  \vskip0.12cm
  \begin{beamercolorbox}[sep=3pt,rounded=true,wd=\linewidth]{block body}
    \centering\small
    \textcolor{scured}{\bfseries MQAR (Multi-Query Associative Recall)}: 检验长上下文中的多查询关联回忆能力.\\[2pt]
    关键不是只读完整个上下文, 而是在大量干扰信息中保持并读出正确的键值绑定.
  \end{beamercolorbox}
  \end{minipage}
\end{frame}

\begin{frame}{全注意力: 记得准, 但代价高}
  \centering
  \begin{minipage}{0.92\textwidth}
  \vspace{0.075cm}
  \small
  \setbeamerfont{block title}{size=\small}
  标准 Transformer 可以让每个 token 直接访问所有历史 token。

  \vskip0.00cm
  \begin{columns}[c,onlytextwidth]
    \begin{column}{0.62\linewidth}
      \centering
      \begin{tikzpicture}
        \node[anchor=south west,inner sep=0] (attentionimg) at (0,0)
          {\makebox[\linewidth][c]{\includegraphics[width=1.08\linewidth,keepaspectratio]{20260523/p04/attention-mechanism-top.jpg}}};
        \begin{scope}[x={(attentionimg.south east)},y={(attentionimg.north west)}]
          \fill[white] (0.935,0.00) rectangle (1.00,0.31);
        \end{scope}
      \end{tikzpicture}
    \end{column}
    \begin{column}{0.34\linewidth}
      \begin{block}{优势}
        \footnotesize
        \vspace{0.08cm}
        直接比较所有历史位置\\
        关联回忆能力强
        \vspace{0.08cm}
      \end{block}

      \vskip0.02cm
      \begin{block}{代价}
        \footnotesize
        \vspace{0.08cm}
        所有 token 两两比较\\
        长序列下成本快速上升
        \vspace{0.08cm}
      \end{block}
    \end{column}
  \end{columns}

  \vskip0.02cm
  {\centering\includegraphics[width=0.66\textwidth,keepaspectratio]{20260523/p04/attention-complexity-table.jpg}\par}
  \end{minipage}
\end{frame}

\begin{frame}{高效模型: 把历史压缩进有限状态}
  \centering
  \begin{minipage}{0.92\textwidth}
  \vspace{0.03cm}
  \small
  \setbeamerfont{block title}{size=\small}
  线性注意力, 状态空间模型, Mamba, GDN 等方法把历史信息压缩进可递推更新的状态.

  \vskip0.08cm
  \begin{columns}[c,onlytextwidth]
    \begin{column}{0.67\linewidth}
      \begin{columns}[T,totalwidth=\linewidth]
        \begin{column}{0.47\linewidth}
          \begin{block}{状态递推与读出}
            \centering
            \vspace{0.04cm}
            {\large\(\displaystyle \symbfit{S}_t = \symbfit{S}_{t-1} + \symbfit{k}_t\symbfit{v}_t^{\symup{T}}\)}\\[4pt]
            {\large\(\displaystyle \symbfit{o}_t = \symbfit{q}_t^{\symup{T}}\symbfit{S}_t\)}
            \vspace{0.04cm}
          \end{block}
        \end{column}

        \begin{column}{0.47\linewidth}
          \begin{block}{共同思想}
            \footnotesize
            \vspace{0.06cm}
            历史不再逐 token 保留\\[2pt]
            通过 \(\symbfit{S}_t\) 递推写入\\[2pt]
            当前 query 从状态读出
            \vspace{0.06cm}
          \end{block}
        \end{column}
      \end{columns}

      \vskip0.20cm
      {\centering\includegraphics[width=0.94\linewidth,keepaspectratio]{20260523/p05/linear-complexity-table.jpg}\par}
    \end{column}

    \begin{column}{0.29\linewidth}
      \centering
      \includegraphics[height=4.45cm,keepaspectratio]{20260523/p05/gla-structure-white.png}
    \end{column}
  \end{columns}
  \end{minipage}
\end{frame}

\begin{frame}{统一视角: 高效模型都在维护可读写的状态}
  \centering
  \begin{minipage}{0.92\textwidth}
  \vspace{0.03cm}
  \small
  不同高效序列模型的递推形式不同, 但大多可以抽象为:
  \textcolor{scured}{状态写入} + \textcolor{scured}{状态读出}.

  \vskip0.08cm
  {\centering\includegraphics[width=0.86\textwidth,keepaspectratio]{20260523/p05/model-recurrence-table.jpg}\par}

  \vskip0.10cm
  \begin{beamercolorbox}[sep=4pt,rounded=true,wd=\linewidth]{block body}
    \centering\small
    \textcolor{scured}{\bfseries 差异在于如何更新状态; 共同点在于都把长历史压缩到可递推的记忆状态中.}
  \end{beamercolorbox}
  \end{minipage}
\end{frame}

\begin{frame}{压缩的代价: 容量不足会损伤检索}
  \centering
  \begin{minipage}{0.92\textwidth}
  \vspace{0.02cm}
  \small
  有限状态压缩把多组键值绑定写入同一个状态, 容量不足时查询读出会发生混叠.

  \vskip0.08cm
  \begin{columns}[T,onlytextwidth]
    \begin{column}{0.50\linewidth}
      \vskip0.16cm
      \thinsectionbar{机制直觉: 聚合会混入无关信息}
      \vskip0.03cm
      {\centering\includegraphics[width=\linewidth,height=4.30cm,keepaspectratio]{20260523/p07/compression-confusion.jpg}\par}
    \end{column}

    \begin{column}{0.46\linewidth}
      \thinsectionbar{性能证据: Retrieval 受损}
      \vskip0.03cm
      {\centering\includegraphics[width=\linewidth,height=2.23cm,keepaspectratio]{20260523/p07/deltanet-retrieval-table.jpg}\par}

      \vskip0.12cm
      \thinsectionbar{容量证据: 维度越大, 记忆越强}
      \vskip0.03cm
      {\centering\includegraphics[width=\linewidth,height=2.35cm,keepaspectratio]{20260523/p07/mqar-capacity.jpg}\par}
    \end{column}
  \end{columns}

  \vskip0.08cm
  \begin{beamercolorbox}[sep=4pt,rounded=true,wd=\linewidth]{block body}
    \centering\small
    \textcolor{scured}{\bfseries 高效模型的问题不是不能压缩, 而是压缩后还能不能准确保存并读出键值绑定.}
  \end{beamercolorbox}
  \end{minipage}
\end{frame}

\begin{frame}[t]{现有高效模型的研究缺口}
  \centering
  \vspace{0.28cm}
  \begin{minipage}{0.92\textwidth}
  \vspace{0.03cm}
  \small
  已有高效长上下文模型分别从遗忘, 改写和索引三条路线缓解状态压缩问题.

  \vskip0.20cm
  {\scriptsize
  \setlength{\tabcolsep}{4pt}
  \renewcommand{\arraystretch}{1.18}
  \arrayrulecolor{scured!35}
  \begin{tabular}{p{0.13\linewidth}p{0.29\linewidth}p{0.24\linewidth}p{0.24\linewidth}}
    \rowcolor{scured}
    \rule[-0.10cm]{0pt}{0.50cm}\textcolor{white}{\bfseries 方法思路} &
    \rule[-0.10cm]{0pt}{0.50cm}\textcolor{white}{\bfseries 代表模型} &
    \rule[-0.10cm]{0pt}{0.50cm}\textcolor{white}{\bfseries 主要解决什么} &
    \rule[-0.10cm]{0pt}{0.50cm}\textcolor{white}{\bfseries 仍然不足} \\
    \hline
    \tablebodycell{\linewidth}{\textcolor{scured}{\bfseries 动态遗忘}} &
    \tablebodycell{\linewidth}{GLA / Mamba} &
    \tablebodycell{\linewidth}{控制历史信息保留,\\避免无关状态长期累积} &
    \tablebodycell{\linewidth}{不直接解决槽内混叠} \\
    \hline
    \tablebodycell{\linewidth}{\textcolor{scured}{\bfseries 定向改写}} &
    \tablebodycell{\linewidth}{DeltaNet / GDN} &
    \tablebodycell{\linewidth}{用输入相关的更新规则\\改写旧记忆} &
    \tablebodycell{\linewidth}{缺少显式分槽组织} \\
    \hline
    \tablebodycell{\linewidth}{\textcolor{scured}{\bfseries 量化索引}} &
    \tablebodycell{\linewidth}{Transformer-VQ / LongVQ} &
    \tablebodycell{\linewidth}{VQ 压缩全局历史,\\线性化长序列注意力} &
    \tablebodycell{\linewidth}{槽内混合误差\\校正不足} \\
    \hline
    \rowcolor{sculight}
    \tablebodycell{\linewidth}{\textcolor{scured}{\bfseries 本课题关注}} &
    \tablebodycell{\linewidth}{\textcolor{scured}{\bfseries Flash-VQRM}} &
    \tablebodycell{\linewidth}{\textcolor{scured}{\bfseries 有限状态混叠\\与槽内读出误差}} &
    \tablebodycell{\linewidth}{\textcolor{scured}{\bfseries 待系统验证}} \\
    \hline
  \end{tabular}
  }

  \vskip0.22cm
  \begin{beamercolorbox}[sep=5pt,rounded=true,wd=\linewidth]{block body}
    \centering\small
    \textcolor{scured}{\bfseries 关键缺口: 分槽之后, 同一记忆槽内的读出误差仍需要进一步校正.}
  \end{beamercolorbox}
  \end{minipage}
\end{frame}

\begin{frame}[t]{我们的思路: 先分槽, 再校正}
  \centering
  \vspace{0.04cm}
  \begin{minipage}{0.92\textwidth}
  \small
  本课题将该结构命名为 \textcolor{scured}{\bfseries Flash-VQRM} (Flash Vector-Quantized Residual Memory Attention): 量化索引先把历史分到记忆槽, 槽内残差再校正粗粒度读出误差.

  \vskip0.08cm
  {\centering
    \includegraphics[width=\linewidth,height=5.55cm,keepaspectratio]{20260523/p09/residual-memory.jpg}\par
  }

  \vskip0.00cm
  \begin{beamercolorbox}[sep=3pt,rounded=true,wd=\linewidth]{block body}
    \centering\small
    \textcolor{scured}{\bfseries 设计直觉: 分槽减少跨类冲突, 残差补偿槽内误差.}
  \end{beamercolorbox}
  \end{minipage}
\end{frame}

\begin{frame}[t]{方法概览: Flash-VQRM}
  \centering
  \vspace{0.06cm}
  \begin{minipage}{0.92\textwidth}
  \small
  \setbeamerfont{block title}{size=\small}
  模型把一次状态读写拆成三部分: \textcolor{scured}{分槽统计}, \textcolor{scured}{残差校正}, 和 \textcolor{scured}{顶层融合}.

  \vskip0.06cm
  \begin{columns}[T,onlytextwidth]
    \begin{column}{0.64\linewidth}
      \centering
      \includegraphics[width=1.14\linewidth,height=5.55cm,keepaspectratio]{20260523/p10/model-structure.jpg}
    \end{column}

    \begin{column}{0.34\linewidth}
      \vskip-0.16cm
      \begin{block}{Coarse memory: 先分槽统计}
        \footnotesize
        \vspace{0.04cm}
        VQ 将相近键路由到同一记忆槽, coarse memory 保存槽级粗粒度历史.
        \vspace{0.04cm}
      \end{block}

      \vskip0.03cm
      \begin{block}{Residual memory: 再校正误差}
        \footnotesize
        \vspace{0.04cm}
        residual memory 记录槽内细节偏差, 用于补偿粗粒度读出的误差.
        \vspace{0.04cm}
      \end{block}

      \vskip0.03cm
      \begin{block}{Aggregation: 融合三路输出}
        \footnotesize
        \vspace{0.04cm}
        顶层聚合把 \(O_{local}\), \(O_{coarse}\), \(O_{res}\) 合成为最终输出 \(O_t\).
        \vspace{0.04cm}
      \end{block}
    \end{column}
  \end{columns}

  \vskip0.04cm
  \begin{beamercolorbox}[sep=3pt,rounded=true,wd=\linewidth]{block body}
    \centering\small
    \textcolor{scured}{\bfseries 关键流程: VQ 分槽组织历史, residual memory 负责槽内误差补偿.}
  \end{beamercolorbox}
  \end{minipage}
\end{frame}

\begin{frame}[t]{MQAR 预研: 131k 容量下高负载更稳}
  \centering
  \vspace{0.06cm}
  \begin{minipage}{0.92\textwidth}
  \small
  \setbeamerfont{block title}{size=\small}
  同为 \textcolor{scured}{\bfseries 131k active capacity}, Flash-VQRM 在 hard slice 1024x256 上保留更高准确率.

  \vskip0.08cm
  \begin{columns}[T,onlytextwidth]
    \begin{column}{0.57\linewidth}
      \centering
      \vskip0.15cm
      \includegraphics[width=1.03\linewidth,height=5.08cm,keepaspectratio]{20260523/p11/mqar-metric-grouped-bar.png}
    \end{column}

    \begin{column}{0.41\linewidth}
      \centering
      \vskip0.02cm
      \makebox[\linewidth][c]{\includegraphics[width=1.12\linewidth,height=4.78cm,keepaspectratio]{20260523/p11/mqar-capacity-fair-table.png}}

      \vskip-0.02cm
      \begin{block}{主要读数 (1024x256)}
        \footnotesize
        \vspace{0.04cm}
        Flash cb64-r16: \textcolor{scured}{\bfseries 0.969}.\\
        GDN h1-ev8: 0.812, GDN h2-ev16: 0.799.
        \vspace{0.04cm}
      \end{block}
    \end{column}
  \end{columns}

  \vskip0.05cm
  \begin{beamercolorbox}[sep=3pt,rounded=true,wd=\linewidth]{block body}
    \centering\small
    \textcolor{scured}{\bfseries 结论: 同等容量下, Flash-VQRM 在高负载 hard slice 上更稳, 说明分槽加校正能缓解状态压缩误差.}
  \end{beamercolorbox}
  \end{minipage}
\end{frame}

\begin{frame}[t]{长度外推预研: Flash-VQRM 衰减更慢}
  \centering
  \vspace{0.00cm}
  \begin{minipage}{0.92\textwidth}
  \small
  \setbeamerfont{block title}{size=\small}
  Longer-MQAR 不重新训练, 直接评估超过训练长度后的关联回忆能力.

  \vskip0.02cm
  \begin{columns}[T,onlytextwidth]
    \begin{column}{0.56\linewidth}
      \centering
      \makebox[\linewidth][c]{\includegraphics[width=1.12\linewidth,height=5.20cm,keepaspectratio]{20260523/p12/longer-mqar-accuracy-curve.pdf}}
    \end{column}

    \begin{column}{0.42\linewidth}
      \centering
      \vskip0.02cm
      \makebox[\linewidth][c]{\includegraphics[width=1.18\linewidth,height=3.75cm,keepaspectratio]{20260523/p12/longer-mqar-accuracy-table.pdf}}

      \vskip-0.28cm
      \makebox[\linewidth][c]{%
      \begin{minipage}{1.18\linewidth}
      \begin{block}{主要读数}
        \footnotesize
        \vspace{0.02cm}
        2048x512: Flash cb64-r16 为 \textcolor{scured}{\bfseries 0.823}, GDN 为 0.355;\\
        4096x1024: Flash cb64-r16 为 \textcolor{scured}{\bfseries 0.469}, GDN 为 0.107.
        \vspace{0.02cm}
      \end{block}
      \end{minipage}}
    \end{column}
  \end{columns}

  \vskip0.32cm
  \begin{beamercolorbox}[sep=3pt,rounded=true,wd=\linewidth]{block body}
    \centering\small
    \textcolor{scured}{\bfseries 结论: 不重训直接外推时, Flash-VQRM 在 2048 到 8190 slice 上衰减更慢, 保留更强关联回忆能力.}
  \end{beamercolorbox}
  \end{minipage}
\end{frame}

\begin{frame}[t]{本课题拟解决的三个关键问题}
  \centering
  \vspace{0.62cm}
  \begin{minipage}{0.92\textwidth}
  \small
  前面预研说明方法方向可行, 后续重点是回答三个机制性问题.

  \vskip0.16cm
  \begingroup
  \footnotesize
  \setlength{\tabcolsep}{5pt}
  \renewcommand{\arraystretch}{1.60}
  \begin{tabular}{>{\raggedright\arraybackslash}p{0.31\linewidth}>{\raggedright\arraybackslash}p{0.29\linewidth}>{\raggedright\arraybackslash}p{0.30\linewidth}}
    \rowcolor{scured}
    \textcolor{white}{\bfseries 关键问题} &
    \textcolor{white}{\bfseries 对应机制} &
    \textcolor{white}{\bfseries 验证方式} \\
    \hline
    \rowcolor{sculight}
    \textcolor{scured}{\bfseries Q1: 如何减少记忆混叠?} &
    量化索引分槽, 让不同历史信息分开组织. &
    MQAR 高负载切片, capacity-fair 强基线对比. \\
    \hline
    \textcolor{scured}{\bfseries Q2: 同槽读出误差如何校正?} &
    residual memory 记录槽内偏差, 补偿粗粒度读出误差. &
    coarse-only / no-residual 消融, cb / rank 变化分析. \\
    \hline
    \rowcolor{sculight}
    \textcolor{scured}{\bfseries Q3: 如何兼顾性能与效率?} &
    限制 active capacity 与状态开销, 提高单位容量收益. &
    准确率, 吞吐和显存统一评估, 结合 Longer-MQAR 分析. \\
  \end{tabular}
  \endgroup

  \vskip0.28cm
  \begin{beamercolorbox}[sep=4pt,rounded=true,wd=\linewidth]{block body}
    \centering\small
    \textcolor{scured}{\bfseries 核心闭环: 先解决有限状态混叠, 再校正槽内误差, 最后评估性能与效率的可用平衡.}
  \end{beamercolorbox}
  \end{minipage}
\end{frame}

\begin{frame}[t]{总体研究路线图}
  \centering
  \vspace{0.04cm}
  \begin{minipage}{0.94\textwidth}
  \small

  \vskip0.02cm
  \centering
  \includegraphics[width=\linewidth,height=5.72cm,keepaspectratio]{20260523/p14/research-roadmap.png}

  \vskip0.05cm
  \begin{beamercolorbox}[sep=3pt,rounded=true,wd=\linewidth]{block body}
    \centering\small
    \textcolor{scured}{\bfseries 路线重点: 围绕混叠抑制, 槽内误差校正和性能效率平衡, 形成 Flash-VQRM 的机制验证闭环.}
  \end{beamercolorbox}
  \end{minipage}
\end{frame}

\begin{frame}[t]{感谢各位老师}
  \centering
  \vfill
  
\includegraphics[height=1.10cm]{common/scu-brand/sichuan-university-logo-cropped.png}

  \vskip0.55cm
  \makebox[\textwidth][c]{%
    \begin{beamercolorbox}[sep=8pt,center,rounded=true,shadow=true,wd=0.66\textwidth]{title}
      \usebeamerfont{title}请各位老师批评指正
    \end{beamercolorbox}%
  }

  \vfill
\end{frame}
